\documentclass[11pt]{letter}
\usepackage[margin=1in]{geometry}
\usepackage{hyperref}
\usepackage{lmodern}
\signature{Chenxi He\\Cavendish Laboratory\\University of Cambridge\\\texttt{ch2067@cam.ac.uk}}
\address{Chenxi He\\Cavendish Laboratory\\University of Cambridge\\Cambridge CB3 0HE, UK}
\date{June 10, 2026}

\begin{document}
\begin{letter}{Editorial Office\\Communications Physics\\Nature Portfolio}
\opening{Dear Editor,}

We are pleased to submit our manuscript ``\emph{Learned variable-support priors
accelerate symbolic regression of physical laws}'' for consideration as an
Article in \emph{Communications Physics}.

\textbf{What we report.} Symbolic regression (SR) is the method of choice for
recovering interpretable physical laws from data, but its cost scales steeply
with the number of candidate variables and operators, most of which never
appear in the true equation. We introduce DenoisedSR, a learned front-end that
reads a small set of observations ($q \approx 100$) and predicts which input
columns are relevant. The front-end couples a graph-attention network trained
over a physical-law co-occurrence graph with a statistical-feature random
forest; its output is a reduced variable/operator set that any conventional SR
backend can consume. DenoisedSR is therefore solver-agnostic by design.

\textbf{Why this matters to a broad physics readership.}
Across four external benchmarks of physical and symbolic laws---AI-Feynman,
Strogatz, Nguyen and the SRSD-Feynman suite of Matsubara et al.---the predictor
attains recall $\ge 0.96$ and never catastrophically discards a true variable.
Supplied to a fixed-budget evolutionary backend (PySR), the prior raises
exact-law recovery from $42\pm13\%$ to a seed-stable $60\pm0\%$, shrinks the
candidate-column set by $85\%$, and speeds time-to-solution by up to
$6\times$. The same qualitative gain reproduces on two structurally different
backends: classical genetic programming (gplearn) and the 2026 Nature
Computational Science cover algorithm PSE (Ruan et al., parallel GPU
enumeration); on PSE we observe a $+34$ to $+67$ percentage-point lift in
exact recovery. The prior is observation-conditioned and adds only
${\sim}40$~ms per task, making it a high-leverage, low-risk preprocessing
step rather than a competing search algorithm. Beyond the headline metric,
our analysis dissects why the prior works (variable selection from
column-target statistics, operator selection from physics-knowledge-graph
priors), establishes that the gain widens with distractor density rather than
benchmark choice, and documents the natural boundary of the approach (the
prior is neutral when distractor density is already low).

\textbf{Fit for the journal.}
This work sits squarely at the intersection of data-driven physics and
machine-learning methodology---two communities \emph{Communications Physics}
serves well. We believe it will be of broad interest to physicists who use
or develop symbolic regression, and to the wider community working on
scientific machine learning. The contribution is complementary to the
dimensional-analysis and physics-aware priors that already strengthen
scientific SR: we contribute the upstream step of deciding \emph{which
measured quantities enter the law at all}.

\textbf{Open data and code.} All experimental artifacts and trained model
weights are publicly released on GitHub
(\url{https://github.com/ChenxiHeCam/DenoisedSR}) and archived on Zenodo
(\href{https://doi.org/10.5281/zenodo.20584889}{10.5281/zenodo.20584889}).
The release includes paper-reproducing evaluation scripts, all recorded
JSON artifacts that back every figure and table, and the trained
graph-attention weights.

\textbf{Originality.} The manuscript reports original work, has not been
published elsewhere, and is not under consideration by any other journal.

\textbf{Suggested reviewers} (with non-overlapping expertise across SR
methodology, neural-symbolic methods, and physical-law discovery; please
substitute or amend as the editor sees fit):

\begin{itemize}\itemsep0pt
\item \textbf{Miles Cranmer}, University of Cambridge --- creator of PySR;
expert in symbolic regression for physics. \emph{[suggest: \texttt{mc2473@cam.ac.uk}]}
\item \textbf{Silviu-Marian Udrescu}, Massachusetts Institute of Technology
--- first author of AI-Feynman; symbolic regression and physics. \emph{[suggest: \texttt{sudrescu@mit.edu}]}
\item \textbf{Yoshitomo Matsubara}, OMRON SINIC X --- author of SRSD-Feynman benchmark; SR for scientific discovery. \emph{[suggest: \texttt{yoshitomo.matsubara@sinicx.com}]}
\item \textbf{Kai Ruan}, Renmin University of China --- first author of PSE
(Nature Computational Science 2026); SR scalability. \emph{[suggest: via group page]}
\item \textbf{Pierre-Alexandre Kamienny}, Meta AI --- end-to-end neural symbolic regression; complementary perspective on assignment problem.
\end{itemize}

\textbf{Opposed reviewers.} None.

We hope you will find the work suitable for review.

\closing{Sincerely yours,}

\end{letter}
\end{document}
